\section*{Bab \ref{ch:b3:learning-memory} --- Plastisitas Saraf, Pembelajaran, dan Ingatan}

\begin{solution}{exo:b3:learning-memory:1}
\omterm{def:b3:learning-memory:kinds}{Ingatan deklaratif} (fakta, kejadian) ditiadakan pada H.M. --- ia tidak dapat
mengingat apa pun yang baru --- sementara ingatan nondeklaratifnya selamat: ia mempelajari
gambar cermin secara normal tanpa mengingat sesinya. Keterampilan:
striatum dan \omterm{def:b3:nervous-systems:plan}{otak kecil}; rasa takut terlazimkan: amigdala; refleks motorik
terlazimkan: \omterm{def:b3:nervous-systems:plan}{otak kecil}; pemrimaan: korteks indra.
\end{solution}

\begin{solution}{exo:b3:learning-memory:2}
Ketika sebuah sel prasinapsis berulang kali menyumbang pada penyalaan sebuah sel
pascasinapsis, sambungan keduanya menguat. Karena itu LTP bersifat khas-masukan (hanya
sinapsis yang aktif, yang ikut serta, yang berubah), asosiatif (sebuah masukan
lemah yang aktif sementara selnya menyala di bawah yang kuat menjadi kuat),
dan kooperatif (cukup banyak masukan harus bekerja bersama untuk menyalakan selnya).
\end{solution}

\begin{solution}{exo:b3:learning-memory:3}
Ia menghantar hanya ketika glutamatnya terikat (sel prasinapsisnya menyala)
\emph{dan} selaputnya terdepolarisasi cukup untuk mengusir magnesiumnya
(sel pascasinapsisnya aktif): kedua syarat \omterm{def:b3:learning-memory:ltp}{kaidah Hebb} di dalam
satu protein. AP5 sebelum tetanusnya: tidak ada kalsium yang masuk, tidak ada LTP, meskipun
penularan lewat \omterm{prop:b3:learning-memory:nmda}{reseptor AMPA} normal. AP5 sesudahnya: LTP yang sudah
terimbas tidak terpengaruh --- reseptornya diperlukan bagi pengimbasannya, bukan
pemeliharaannya.
\end{solution}

\begin{solution}{exo:b3:learning-memory:4}
Jangka pendek: perubahan kovalen pada protein yang ada --- PKA
memfosforilasi lalu menutup sebuah saluran kalium, yang melebarkan lecutannya
dan menambah pelepasannya; tanpa protein baru. Jangka panjang: PKA mengaktifkan CREB,
gen baru ditranskripsi dan terminal sinapsis baru tumbuh. Sebuah
penghambat sintesis protein yang diberikan selama latihan yang berulang membiarkan
fasilitasi jangka pendeknya utuh lalu meniadakan ingatan yang diukur berhari-hari
kemudian; bila diberikan sesudahnya, ia tidak melakukan apa pun.
\end{solution}

\begin{solution}{exo:b3:learning-memory:5}
$0.7^{n} \le 0.1$: $n \ge 6.5$, sehingga $7$ percobaan. Dari $0.9\lambda$,
$0.9\times 0.7^{n} < 0.1$: sekali lagi $7$ percobaan. Setangkup karena $\alpha$ yang
sama mengatur kedua arahnya. Pada hewan, pemunahannya biasanya
lebih lambat dan ingatan aslinya kembali sesudah sebuah istirahat (pemulihan
spontan): pemunahan adalah pembelajaran baru yang menekan yang lama, bukan penghapusannya.
\end{solution}

\begin{solution}{exo:b3:learning-memory:6}
Vektor eigennya $(1,1)/\sqrt{2}$ bernilai eigen $1.5$ dan $(1,-1)/\sqrt{2}$
dengan $0.5$; Hebb memutar $\mathbf{w}$ ke arah $(1,1)$, dan neuronnya berakhir
mendeteksi kedua masukannya menyala bersama. Langkah Oja: $y = 1$,
$\Delta\mathbf{w} = 0.1\times 1\times((1,1) - 1\times(1,0)) = (0,0.1)$,
sehingga $\mathbf{w} = (1, 0.1)$.
\end{solution}

\begin{solution}{exo:b3:learning-memory:7}
$500\times 0.05 = 25$ ion bebas; $25/6\times 10^{23} = 4.2\times 10^{-23}$
mol di dalam $10^{-16}$ L: \qty{0.42}{\micro mol/L}, yakni empat kali aras
istirahatnya tetapi di bawah $K_{d}$ kalmodulinnya --- pengaktifan sebagian saja;
beberapa pembukaan saluran bersama-sama diperlukan untuk mencapai jangkauan mikromolar
yang menyalakan CaMKII.
\end{solution}

\begin{solution}{exo:b3:learning-memory:8}
A dilatih sampai $V_{A} = \lambda$: dopaminnya meledak pada percobaan
awalnya lalu membisu ketika ganjarannya menjadi teramalkan. Pada percobaan
majemuknya, ramalan $V_{A} + V_{B} = \lambda$ sudah tepat,
galatnya nol, dopaminnya bisu, dan $V_{B}$ tidak pernah berubah: B
terblokir. Ketika diuji sendirian, B tidak memicu tanggapan; bila sebuah ganjaran lalu mengikuti B
sendirian, ganjaran itu tak terduga dan dopaminnya meledak.
\end{solution}

\begin{solution}{exo:b3:learning-memory:9}
(a) Tidak ada LTP di \omterm{def:b3:learning-memory:hippocampus}{CA1} dan sebuah kekurangan ingatan ruang, dengan pembelajaran lain
utuh. (b) Ingatan jangka pendeknya normal, ingatan jangka panjangnya tidak ada pada
\qty{24}{h}. (c) Tidak berpengaruh --- \omterm{def:b3:learning-memory:kinds}{pemantapannya} sudah selesai (kecuali bila ingatannya
diaktifkan kembali, ketika ia menjadi labil lagi). (d) Mencitnya membeku di dalam
konteks yang aman: ingatannya dipanggil secara buatan dengan mengaktifkan kembali
sel yang menyandinya.
\end{solution}

\begin{solution}{exo:b3:learning-memory:10}
$\Delta|\mathbf{w}|^{2} \approx 2\mathbf{w}\cdot\Delta\mathbf{w} = 2\eta
y^{2} \ge 0$: panjangnya tidak pernah menurun lalu tumbuh kapan pun neuronnya
menyala. Pada sebuah titik tetap Oja, $C\mathbf{w} = (\mathbf{w}^{\mathsf{T}}C
\mathbf{w})\mathbf{w}$, sehingga $\mathbf{w}$ sebuah vektor eigen dengan $\lambda =
\lambda|\mathbf{w}|^{2}$, sehingga $|\mathbf{w}| = 1$. Pertumbuhan tanpa batas
akan menjenuhkan tiap sinapsis, meniadakan keterpilihannya, dan mendorong
pemacuan yang lepas kendali; penormal biologisnya adalah penskalaan sinapsis (semua
sinapsis sebuah neuron diskalakan turun ketika ia terlalu aktif), LTD, dan
penskalaan turun sinapsis selama tidur.
\end{solution}

\begin{solution}{exo:b3:learning-memory:11}
Tikus: $0.14\times 3\times 10^{5} \approx 4\times 10^{4}$; manusia: $0.14
\times 2\times 10^{6} \approx 3\times 10^{5}$. Taksirannya menganggap
pola acak yang padat dan kesambungan penuh; pola nyatanya jarang
(yang menaikkan kapasitasnya banyak) dan terstruktur, dan \omterm{def:b3:learning-memory:kinds}{hipokampusnya} adalah
sebuah simpanan sementara yang menyerahkan ingatan kepada sebuah korteks berneuron $10^{10}$
--- simpanan seumur hidupnya adalah korteksnya, dan angka \omterm{def:b3:learning-memory:kinds}{hipokampusnya} adalah sebuah
ukuran penyangga, bukan sebuah cacah ingatan.
\end{solution}

\begin{solution}{exo:b3:learning-memory:12}
Masukan mata yang terbuka berkorelasi dengan penyalaan sel korteksnya
lalu diperkuat; masukan mata yang tertutup tidak, dan di bawah
persaingan bagi sebuah bobot total yang terbatas mereka melemah --- yakni Hebb ditambah
penormalan. Dengan kedua matanya tertutup, tidak ada masukan yang dipihak, sehingga
keduanya menyimpan wilayah (yang lemah) dan pengaruhnya lebih ringan. Masanya menutup
ketika sirkuit penghambatnya matang, jala perineuron mengaku di sekeliling
sinapsisnya, dan subunit \omterm{prop:b3:learning-memory:nmda}{reseptor NMDA-nya} berubah menjadi sebuah
bentuk yang kurang mengizinkan.
\end{solution}

\begin{solution}{pb:b3:learning-memory:1}
\textbf{1.} $10\times 10^{3} = 10^{4}$ ion; \qty{2}{\%} bebas, yakni $200$;
$200/6\times 10^{23} = 3.3\times 10^{-22}$ mol di dalam $8\times 10^{-17}$ L:
\qty{4.2}{\micro mol/L}.
\textbf{2.} Ya, di atas ambang \qty{2}{\micro mol/L}. Dua
reseptor: $40$ ion bebas, \qty{0.83}{\micro mol/L} --- di bawahnya.
\textbf{3.} Pada \qty{0.83}{\micro mol/L} kalsineurin aktif dan CaMKII
tidak: \omterm{prop:b3:learning-memory:nmda}{reseptor AMPA} disingkirkan dan sinapsisnya melemah (LTD). Satu
ion, dua enzim berafinitas berbeda: sebuah kenaikan yang besar dan cepat mencapai
kinasenya, sebuah kenaikan yang kecil dan lambat hanya fosfatasenya.
\textbf{4.} Sebuah masukan kedua dalam beberapa tetapan waktu menambahkan kalsiumnya
pada sisa yang pertama; sesudah \qty{50}{ms} hampir tidak ada yang
tersisa. Tuntutan kebetulan NMDA dan peluruhan \qty{20}{ms}
bersama-sama memberi sebuah jendela sekitar \qty{20}{ms} pada tiap sisinya --- yakni
jendela saat lecutannya.
\textbf{5.} Penyangganya menangkap kebanyakan kalsiumnya dalam sepersekian
mikrometer dan leher duri yang sempit memperlambat pembauran ke dendritnya,
sehingga kenaikannya tetap di dalam duri yang aktif; sebuah tetangga yang
berjarak \qty{0.5}{\micro m} tidak melihat apa-apa, dan hanya sinapsis yang aktif yang
berubah.
\textbf{6.} $(70 - 30)/5 = 8$ masukan serentak, atau sebuah potensial aksi
yang merambat balik dari badan selnya: sebuah sinapsis tunggal tidak dapat membuka halangan
\omterm{prop:b3:learning-memory:nmda}{reseptor NMDA-nya} sendiri --- \omterm{def:b3:structural-biology:allostery}{kekooperatifannya} terbangun ke dalam fisikanya.
\textbf{7.} $(1,1)/\sqrt{2}$ bernilai eigen $1.6$; $(1,-1)/\sqrt{2}$
dengan $0.4$.
\textbf{8.} $C\mathbf{w}_{0} = (1, 0.6)$, $\mathbf{w}_{1} = (1.10,
0.06)$, sudut terhadap $(1,1)$: $45^{\circ} - 3.1^{\circ} = 41.9^{\circ}$.
$\mathbf{w}_{2} = (1.21, 0.13)$: $38.8^{\circ}$. $\mathbf{w}_{3} = (1.34,
0.22)$: $35.8^{\circ}$.
\textbf{9.} $w_{2}/w_{1} \to 1$ sementara $|\mathbf{w}|$ tumbuh tanpa batas
(dengan faktor $1.16$ tiap langkah).
\textbf{10.} $C\mathbf{w} = (1.28, 1.28)$; $\mathbf{w}^{\mathsf{T}}C
\mathbf{w} = 2.05$; $\Delta\mathbf{w} = 0.1\times(1.28 - 2.05\times 0.8)
= -0.036$ masing-masing: $\mathbf{w} = (0.76, 0.76)$, yang menyusut ke arah vektor
satuan $(0.71, 0.71)$.
\textbf{11.} Bersama $y = 1.42$; masukan 2 sendirian $y = 0.71$, yakni separuhnya.
\textbf{12.} Masukan lemahnya menyala sementara yang kuat menggerakkan selnya,
sehingga ia berkorelasi dengan keluarannya dan \omterm{def:b3:learning-memory:ltp}{kaidah Hebb} memperkuatnya ---
neuronnya mempelajari kemunculan bersamanya.
\textbf{13.} $1 - 0.85^{n}$: $0.56$, $0.80$, $0.96$.
\textbf{14.} $0.85^{n} \le 0.05$: $n \ge 18.5$, sehingga $19$ percobaan.
\textbf{15.} Tidak: bagian asimtot yang tercapai sesudah $n$ percobaan tidak
bergantung pada $\lambda$; asimtotnya berlipat dua. Sebuah ganjaran yang lebih besar dapat
menaikkan $\alpha$ (kemenonjolan), dan sebuah ambang tingkah laku pada nilai
\emph{mutlaknya} tercapai lebih cepat.
\textbf{16.} $V_{B} = 0$: terblokir. Majemuknya teramalkan sepenuhnya,
galatnya nol, neuron dopaminnya bisu.
\textbf{17.} B sendirian tidak meramalkan apa pun, sehingga galatnya adalah $\lambda$ dan
dopaminnya meledak; $V_{B}$ lalu naik seperti pada perolehannya, $\lambda(1 -
0.85^{n})$.
\textbf{18.} Obatnya menghantarkan ledakan yang berarti ``lebih baik daripada yang
teramalkan'' sesudah apa pun yang mendahuluinya; kaidahnya menaikkan nilai
isyarat dan tindakan itu pada tiap kesempatan, sehingga keduanya memperoleh nilai
lepas dari hasil nyata mana pun --- pencariannya menjadi memaksa.
\textbf{19.} $0.14\times 3\times 10^{5} \approx 4.2\times 10^{4}$.
\textbf{20.} $4.2\times 10^{4}/20 = 2100$ hari, yakni sekitar enam tahun;
\omterm{def:b3:learning-memory:kinds}{pemantapannya} harus memindahkan ingatan ke korteksnya dan pelupaannya harus
membebaskan simpanannya.
\textbf{21.} Menulis pola baru dengan cepat ke dalam korteksnya akan menimpa
sinapsis yang menyandi pola lama (gangguan bencana); \omterm{def:b3:learning-memory:hippocampus}{pemutaran ulang} yang
lambat dan berselang-seling membuat korteksnya memadukan yang baru dengan yang lama,
sementara \omterm{def:b3:learning-memory:kinds}{hipokampusnya} memegang ingatan barunya sementara itu --- yakni sebuah penyangga cepat
yang memberi makan sebuah simpanan lambat.
\textbf{22.} $0.03\times 3\times 10^{5} = 9000$ sel tiap tempat. Pola
yang jarang sedikit saja bertindih, sehingga jauh lebih banyak yang dapat disimpan sebelum
gangguannya merusak pemanggilannya; kapasitasnya naik lebih cepat daripada $N$ ketika
polanya menjadi lebih jarang, dengan ongkos tiap polanya memakai lebih banyak sel
yang bersifat khas.
\textbf{23.} Pemampatan sebesar $20$: sel yang menyala berjarak \qty{400}{ms} pada
larinya menyala berjarak \qty{20}{ms} pada \omterm{def:b3:learning-memory:hippocampus}{pemutaran ulangnya} --- di dalam jendela saat
lecutannya. Maksud pemampatannya adalah membawa urutan yang terlalu lambat untuk
diasosiasikan pada waktu nyata ke dalam jangkauan Hebb.
\textbf{24.} Pada sehari, ingatannya masih bergantung pada \omterm{def:b3:learning-memory:kinds}{hipokampusnya}; pada
sebulan, ia telah dimantapkan ke dalam korteksnya lalu tidak lagi memerlukannya
--- tepat seperti kehilangan H.M. yang berjenjang menurut waktu, dengan ingatan baru yang hilang dan
ingatan jauh yang terluput.
\textbf{25.} Sekitar \qty{4.2}{\micro mol/L} kalsium bebas sesudah sepuluh
pembukaan; $19$ percobaan sampai \qty{95}{\%}; sekitar $4\times 10^{4}$ pola
di dalam \omterm{def:b3:learning-memory:hippocampus}{CA3} tikusnya.
\end{solution}
